\documentclass{beamer}
\usepackage[utf8]{inputenc}

\usetheme{Madrid}
\usecolortheme{default}
\usepackage{amsmath,amssymb,amsfonts,amsthm}
\usepackage{txfonts}
\usepackage{tkz-euclide}
\usepackage{listings}
\usepackage{adjustbox}
\usepackage{array}
\usepackage{tabularx}
\usepackage{gvv}
\usepackage{lmodern}
\usepackage{circuitikz}
\usepackage{tikz}
\usepackage{graphicx}
\usepackage{multicol}

\setbeamertemplate{page number in head/foot}[totalframenumber]

\usepackage{tcolorbox}
\tcbuselibrary{minted,breakable,xparse,skins}

\definecolor{bg}{gray}{0.95}
\DeclareTCBListing{mintedbox}{O{}m!O{}}{%
  breakable=true,
  listing engine=minted,
  listing only,
  minted language=#2,
  minted style=default,
  minted options={%
    linenos,
    gobble=0,
    breaklines=true,
    breakafter=,,
    fontsize=\small,
    numbersep=8pt,
    #1},
  boxsep=0pt,
  left skip=0pt,
  right skip=0pt,
  left=25pt,
  right=0pt,
  top=3pt,
  bottom=3pt,
  arc=5pt,
  leftrule=0pt,
  rightrule=0pt,
  bottomrule=2pt,
  toprule=2pt,
  colback=bg,
  colframe=orange!70,
  enhanced,
  overlay={%
    \begin{tcbclipinterior}
    \fill[orange!20!white] (frame.south west) rectangle ([xshift=20pt]frame.north west);
    \end{tcbclipinterior}},
  #3,
}
\lstset{
    language=C,
    basicstyle=\ttfamily\small,
    keywordstyle=\color{blue},
    stringstyle=\color{orange},
    commentstyle=\color{green!60!black},
    numbers=left,
    numberstyle=\tiny\color{gray},
    breaklines=true,
    showstringspaces=false,
}

\numberwithin{equation}{section}
\lstset{
  language=Python,
  basicstyle=\ttfamily\small,
  keywordstyle=\color{blue},
  stringstyle=\color{orange},
  numbers=left,
  numberstyle=\tiny\color{gray},
  breaklines=true,
  showstringspaces=false
}


\title{Problem 8.3.10}
\author{Sarvesh Tamgade}

\date{\today} 
\begin{document}

\begin{frame}
\titlepage
\end{frame}

\section{Question}
\begin{frame}{Question}
\textbf{Question}:
If the eccentricity of an ellipse is $\frac{5}{8}$ and the distance between its foci is $10$, then find the latus rectum of the ellipse.



\end{frame}

\section{Solution}
\begin{frame}[fragile]
    \frametitle{Solution}
Given:
\begin{itemize}
\item Eccentricity: $e = \frac{5}{8}$
\item Distance between foci: $D = 10$
\end{itemize}

\textbf{Step 1: Express the foci for standard ellipse}

From B.2.7.3, the foci for standard ellipse are at:

\begin{align}
\vec{F} = \pm\sqrt{\frac{f_0}{\lambda_2(1-e^2)}} \vec{e_1}
\end{align}

where $f_0$ is the scalar term in the conic equation and $\lambda_2$ is an eigenvalue of matrix $\vec{V}$.

\textbf{Step 2: Distance between foci}

The distance between foci is:

\begin{align}
D = 2\sqrt{\frac{f_0}{\lambda_2(1-e^2)}}
\end{align}


\end{frame}
\begin{frame}[fragile]
    \frametitle{Solution}
\textbf{Step 3: Solve for $f_0$}

From equation (2):

\begin{align}
\frac{D}{2} &= \sqrt{\frac{f_0}{\lambda_2(1-e^2)}}
\end{align}

Squaring both sides:

\begin{align}
\frac{D^2}{4} &= \frac{f_0}{\lambda_2(1-e^2)}
\end{align}

Therefore:

\begin{align}
f_0 &= \frac{D^2 \lambda_2(1-e^2)}{4}
\end{align}

\end{frame}
\begin{frame}[fragile]
    \frametitle{Solution}
G\textbf{Step 4: Eigenvalues for standard ellipse}

For the standard ellipse, from the conic section theory:

\begin{align}
\lambda_1 &= \frac{||\vec{n}||^2}{1-e^2}
\end{align}

\begin{align}
\lambda_2 &= ||\vec{n}||^2
\end{align}

For the standard form with $||\vec{n}||^2 = 1$:

\begin{align}
\lambda_1 &= \frac{1}{1-e^2}
\end{align}

\begin{align}
\lambda_2 &= 1
\end{align}

\end{frame}

\begin{frame}[fragile]
    \frametitle{Solution}
\textbf{Step 5: Calculate $f_0$ with $\lambda_2 = 1$}

Substituting $\lambda_2 = 1$ into equation (5):

\begin{align}
f_0 &= \frac{D^2(1-e^2)}{4}
\end{align}

Therefore:

\begin{align}
\sqrt{f_0} &= \frac{D\sqrt{1-e^2}}{2}
\end{align}

\textbf{Step 6: Matrix formulation for latus rectum}
The length of latus rectum is given by :

\begin{align}
\ell = \frac{2\sqrt{f_0 \lambda_1}}{\lambda_2}
\end{align}


\end{frame}

\begin{frame}[fragile]
    \frametitle{Solution}
Substituting the values of $f_0$, $\lambda_1$ and $\lambda_2$ 
\begin{align}
\ell &= \frac{D(1-e^2)}{e}
\end{align}

\textbf{Step 7: Substitute the given values}

Given $e = \frac{5}{8}$ and $D = 10$:

\begin{align}
1 - e^2 &= 1 - \left(\frac{5}{8}\right)^2 = 1 - \frac{25}{64} = \frac{39}{64}
\end{align}

Substituting into the latus rectum formula:

\begin{align}
\ell &= \frac{390 \times 8}{64 \times 5} = \frac{3120}{320} = \frac{39}{4}
\end{align}
\end{frame}
\begin{frame}[fragile]
    \frametitle{Solution}
\textbf{Final Answer:}
The length of latus rectum is :

\begin{align}
\boxed{\ell = \frac{39}{4}}
\end{align}
\end{frame}



\section{Graph}
\begin{frame}
    \frametitle{Graph}
    \begin{figure}[htbp]
    \centering
    \includegraphics[width=0.65\linewidth]{FIG/fig1.png}
    \caption{Vector Representation}
    \label{fig:FIG/fig1.png}
\end{figure}
\end{frame}
\section{ C Code}
\begin{frame}[fragile]
\frametitle{C Code }
\begin{lstlisting}[language=C]
#include <math.h>
#include "matfun.h"

double semi_major ( double e , double D ){
    return D / ( 2.0 * e );
}

double semi_minor ( double e , double D ){
    double a = semi_major ( e , D );
    return a * sqrt ( 1.0 - e * e );
}

double focal_distance ( double e , double D ){
    return semi_major ( e , D ) * e;
}
\end{lstlisting}
\end{frame}
\begin{frame}[fragile]
\frametitle{C Code }
\begin{lstlisting}[language=C]
double latus_rectum ( double e , double D ){
    return D * ( 1.0 - e * e ) / e;
}

void ellipse_points(double a , double b ,int    N ,double *x ,double *y){
    for ( int i = 0 ; i < N ; ++i )
    {
        double t = 2.0 * M_PI * i / N;
        x [ i ] = a * cos ( t );
        y [ i ] = b * sin ( t );
    }
}
\end{lstlisting}
\end{frame}
\begin{frame}[fragile]
\frametitle{Python + C Code }
\begin{lstlisting}[language=Python]
import ctypes as ct
import numpy as np
import matplotlib.pyplot as plt
from   numpy.ctypeslib import ndpointer


lib = ct.CDLL ( './libmatfun.so' )

for f in ( 'semi_major' , 'semi_minor' ,
           'focal_distance' , 'latus_rectum' ):
    getattr ( lib , f ).argtypes = [ ct.c_double , ct.c_double ]
    getattr ( lib , f ).restype  =  ct.c_double

lib.ellipse_points.argtypes = [
    ct.c_double , ct.c_double , ct.c_int ,
    ndpointer ( ct.c_double , flags = 'C_CONTIGUOUS' ),
    ndpointer ( ct.c_double , flags = 'C_CONTIGUOUS' )
]
\end{lstlisting}

\end{frame}
\begin{frame}[fragile]
    \frametitle{Python + C Code}
    \begin{lstlisting}
e , D = 5 / 8 , 10

a = lib.semi_major     ( e , D )
b = lib.semi_minor     ( e , D )
c = lib.focal_distance ( e , D )
l = lib.latus_rectum   ( e , D )

print ( l )                           # prints 9.75


N = 1000
x = np.empty ( N )
y = np.empty ( N )

lib.ellipse_points ( a , b , N , x , y )

fig , ax = plt.subplots ( figsize = ( 12 , 8 ) )
\end{lstlisting}
\end{frame}

\begin{frame}[fragile]
    \frametitle{Python + C Code }
    \begin{lstlisting}

ax.plot    ( x , y , 'b-' , lw = 2 )
ax.scatter ( [  c , -c ] , [ 0 , 0 ] , color = 'r' , s = 50 )
ax.plot    ( [  c , -c ] , [ 0 , 0 ] , 'r--' , lw = 1 )

ax.plot    ( [  c ,  c ] , [ -l / 2 ,  l / 2 ] , 'g' , lw = 3 )
ax.plot    ( [ -c , -c ] , [ -l / 2 ,  l / 2 ] , 'g' , lw = 3 )

ax.plot ( [ -a ,  a ] , [ 0 , 0 ] , 'k--' , alpha = .5 )
ax.plot ( [  0 ,  0 ] , [ -b , b ] , 'k--' , alpha = .5 )


\end{lstlisting}
\end{frame}
\begin{frame}[fragile]
    \frametitle{Python + C Code }
    \begin{lstlisting}
ax.text (  a + .3 , 0       , f'A ( {a} )' )
ax.text ( -a - 1 , 0        , f"A' ( -{a} )" )
ax.text ( .3 ,  b + .3      , f'B ( {b:.2f} )' )
ax.text ( .3 , -b - .6      , f"B' ( -{b:.2f} )" )
ax.text (  c + .2 ,  l / 2  , f'ℓ = {l:.2f}' )
ax.text ( -c - 1 ,  l / 2   , f'ℓ = {l:.2f}' )

ax.set_aspect ( 'equal' )
ax.set_xlim   ( -10 , 10 )
ax.set_ylim   ( -8  , 8  )
ax.set_title  ( f'Ellipse  e = {e}   D = {D}' )
ax.grid       ( alpha = .3 )

plt.tight_layout ( )
plt.show ( )



\end{lstlisting}
\end{frame}

\begin{frame}[fragile]
\frametitle{Python Code }
\begin{lstlisting}[language=Python]
import numpy as np
import matplotlib.pyplot as plt

e, D = 5 / 8, 10
l = D * (1 - e ** 2) / e
a = D / (2 * e)
b = a * np.sqrt(1 - e ** 2)
c = a * e
print("The length of latus rectum is: ",l)
t = np.linspace(0, 2 * np.pi, 1000)
x, y = a * np.cos(t), b * np.sin(t)

fig, ax = plt.subplots(figsize=(12, 8))
ax.plot(x, y, 'b-', lw=2, label='Ellipse')

\end{lstlisting}

\end{frame}
\begin{frame}[fragile]
    \frametitle{Python Code}
    \begin{lstlisting}
ax.scatter([c, -c], [0, 0], color='r', s=50, label='Foci $F_1, F_2$')
ax.text(c + .2, 0.3, f'$F_1({c}, 0)$')
ax.text(-c - .2, 0.3, f'$F_2({-c}, 0)$')

ax.plot([-a, a], [0, 0], 'k--', alpha=.5)
ax.plot([0, 0], [-b, b], 'k--', alpha=.5)

ax.text(a + .3, 0, f'$A({a}, 0)$')
ax.text(-a - .9, 0, f'$A\'(-{a}, 0)$')
ax.text(0.3, b + .3, f'$B(0, {b:.2f})$')
ax.text(0.3, -b - .5, f'$B\'(0, -{b:.2f})$')

\end{lstlisting}
\end{frame}

\begin{frame}[fragile]
    \frametitle{Python Code}
    \begin{lstlisting}

ax.plot([c, c], [-l / 2, l / 2], 'g', lw=3, label='Latus rectum')
ax.plot([-c, -c], [-l / 2, l / 2], 'g', lw=3)
ax.text(c + .2, l / 2, f'$\\ell = {l:.2f}$')
ax.text(-c - 1.0, l / 2, f'$\\ell = {l:.2f}$')

ax.set_aspect('equal')
ax.set_xlim(-10, 10)
ax.set_ylim(-8, 8)
ax.set_title(f'Ellipse $e={e}$,  $D={D}$\n$\\ell = {l} = 39/4$')
ax.grid(alpha=.3)
ax.legend()
plt.tight_layout()
plt.savefig("fig1.png")
plt.show()

\end{lstlisting}
\end{frame}
\end{document}
